\chapter{Orbital-space comparison (menu 13)}
\label{ch:menu13}
\menulabel{13}

\section{Purpose}

Two zero-external-reference checks of an orbital space, both read from the
singular values of the overlap matrix between two sets of orbitals,
$M = C_A^{\mathsf{T}} S\, C_B$ (each set orthonormalised internally first; the
source states the result is independent of the orthonormalisation scheme):

\begin{itemize}
  \item the \textbf{subspace fraction}
    $\sigma_F = \lVert M \rVert_F / \sqrt{\min(|A|, |B|)}$ answers ``how much of
    space $B$ is contained in space $A$''. Its intended use is ranking a
    recommended active space against a fuller reference space: the source
    reports that the ordering of the deficit $1 - \sigma_F$ tracks the ordering
    of the energy error. A large $\sigma_F$ cannot exclude that $A$ is too
    large -- the check is one-sided by construction and the report says so;
  \item the \textbf{space-change SVD} compares two sets of the same size (the
    initial and the final active space of one orbital optimisation). Singular
    values near $1$ mean the space hardly moved; a singular value near $0$
    means one initial active orbital was replaced by an unrelated one, i.e. the
    initial space lacked an element needed for the correlated wave function.
\end{itemize}

Both need only two coefficient blocks and the AO overlap matrix of one basis --
exactly what two \code{orca\_2json} exports of the same system provide.

\section{What you need}

\begin{itemize}
  \item two \code{orca\_2json} exports of the same system in the same basis
    (the defaults carry the MO coefficients and the $S$ matrix);
  \item for the identity case -- the same gbw before and after
    \code{orca\_loc} -- which exercises the check's blind spot: a rotation
    \emph{inside} one space must come out as ``unchanged''.
\end{itemize}

\section{Walkthrough}

\lstinputlisting[style=fbkcode, firstline=17, lastline=33,
  title={The menu-13 example session (the N\textsubscript{2} CAS(6,6) exports
  of the same gbw before and after \code{orca\_loc})}]
  {examples/expected/m13_orbital_space.txt}

Each window is given as \code{first last} in ORCA's 0-based orbital numbering,
the same convention as menu~\menu{12}; pressing Enter on a window covers every
orbital of that export.

\section{What you get}

The singular values (descending), the fraction $\sigma_F$ and its deficit, the
smallest singular value, and the two readings with their provisional bands --
along with a report file (\file{<first export>.fbk.md}) that carries the
citations.

\section{How to read the numbers}

\begin{readbox}
\begin{itemize}
  \item \textbf{The identity case is the sanity anchor.} The example compares
    one space with itself seen through a localisation: all six singular values
    are $1.0000$ to better than $10^{-10}$, $\sigma_F = 1$, and the space-change
    reading says ``essentially the same space''. Any rotation inside a space
    must be invisible; a value below $1$ here would mean the two exports do not
    belong to the same basis or system.
  \item \textbf{The bands are provisional and ordered, not absolute.} They are
    anchored on the sources' own numbers ($1 - \sigma_F = 2.9\times10^{-3}$ for
    a deliberately loose start and $6.8\times10^{-5}$ after a localisation in
    the embedding source; a $0.65$--$0.99$ singular-value range in the AVAS
    source). The intended use is comparing several candidate spaces of one
    system; the absolute bands only separate the numerical level from a real
    deficiency.
  \item \textbf{Read the deficit for containment, the smallest singular value
    for a replaced orbital.} A missing direction moves the smallest singular
    value long before it moves the RMS fraction; the report prints both.
  \item \textbf{What the check cannot say.} It cannot tell whether a space is
    too \emph{large} (containment saturates at $1$), and it says nothing about
    the quality of either space -- it is a comparison, not an energy.
\end{itemize}
\end{readbox}

\section{Worked example}

The session above is replayable from the shipped fixtures: the two N2 export
files of the CAS(6,6) chain, copied into the example workspace as
\file{canonical.json} and \file{localized.json}. To use it on your own system,
export the two gbw files with \code{orca\_2json} and give the windows; for the
``recommended versus reference space'' reading, export the gbw of the
calculation that defines the reference space and compare the windows.

\section{Caveats, refusals and related sections}

\begin{refusalbox}
\begin{itemize}
  \item An export without an $S$ matrix is refused (the overlap of the two
    spaces cannot be formed); re-export with the defaults of \code{orca\_2json}.
  \item Two exports of different AO dimensions or different overlap matrices
    are refused as not sharing one basis.
  \item A linearly dependent window is refused: the comparison is defined on
    linearly independent sets.
  \item A window that is not orthonormal in the AO metric is orthonormalised
    first, and the report records that this was necessary (the reading itself is
    independent of the orthonormalisation scheme).
\end{itemize}
\end{refusalbox}

Related: \menu{12} (the exact active-space analysis of the same exports),
\menu{1} (the composition checks that decide which orbitals are in a window),
and Chapter~\ref{ch:criteria} for the threshold table.
